\documentclass{article}
\usepackage{graphicx}
\usepackage{amsmath}

\title{Planetary Gear System N=3}
\author{Rudhdam Shende}
\date{August 2026}

\begin{document}

\maketitle

\section{Calculations}

\subsection{\textbf{Tooth Count Compatibility}}

For a simple planetary gear system:

\begin{equation}
N_r = N_s + 2N_p
\end{equation}

\begin{equation}
N_r = 20 + 2(22) = 64
\end{equation}

Therefore,

\begin{equation}
N_s = 20,\qquad N_p = 22,\qquad N_r = 64
\end{equation}

For three planets, the assembly condition is:

\begin{equation}
\frac{N_s+N_r}{N}
=
\frac{20+64}{3}
=
28
\end{equation}

Hence, the three-planet arrangement is compatible.

\subsection{\textbf{Sun Gear ($N_s$ = 20 Teeth)}}

$m = 2$ mm

\textbf{Pitch Diameter}

\begin{equation}
d_s = mN_s
\end{equation}

\begin{equation}
d_s = 2(20) = 40.000\ \mathrm{mm}
\end{equation}

\begin{equation}
r_s = 20.000\ \mathrm{mm}
\end{equation}

\textbf{Addendum}

\begin{equation}
h_a = m
\end{equation}

\begin{equation}
h_a = 2.000\ \mathrm{mm}
\end{equation}

\textbf{Dedendum}

\begin{equation}
h_f = 1.25m
\end{equation}

\begin{equation}
h_f = 2.500\ \mathrm{mm}
\end{equation}

\textbf{Outside Diameter}

\begin{equation}
d_{a,s} = m(N_s+2)
\end{equation}

\begin{equation}
d_{a,s} = 44.000\ \mathrm{mm}
\end{equation}

\textbf{Root Diameter}

\begin{equation}
d_{f,s} = m(N_s-2.5)
\end{equation}

\begin{equation}
d_{f,s} = 35.000\ \mathrm{mm}
\end{equation}

\textbf{Base Diameter}

\begin{equation}
d_{b,s} = d_s\cos(20^\circ)
\end{equation}

\begin{equation}
d_{b,s} = 37.588\ \mathrm{mm}
\end{equation}

\begin{equation}
r_{b,s} = 18.794\ \mathrm{mm}
\end{equation}

\textbf{TOOTH THICKNESS AT PITCH CIRCLE}

\begin{equation}
t = s = \frac{\pi m}{2}
\end{equation}

\begin{equation}
t = 3.1416\ \mathrm{mm}
\end{equation}

\textbf{CIRCULAR PITCH}

\begin{equation}
p = \pi m
\end{equation}

\begin{equation}
p = 6.283\ \mathrm{mm}
\end{equation}

\textbf{ANGULAR PITCH}

\begin{equation}
\theta = \frac{360^\circ}{N_s}
\end{equation}

\begin{equation}
\theta = 18^\circ
\end{equation}


\subsection{\textbf{Planet Gear ($N_p$ = 22 Teeth)}}

\begin{equation}
N_p = 22
\end{equation}

\textbf{PITCH DIAMETER}

\begin{equation}
d_p = mN_p
\end{equation}

\begin{equation}
d_p = 44.000\ \mathrm{mm}
\end{equation}

\begin{equation}
r_p = 22.000\ \mathrm{mm}
\end{equation}

\textbf{ADDENDUM / DEDENDUM}

\begin{align}
h_a &= 2.000\ \mathrm{mm} \quad \text{(Addendum)}\\
h_f &= 2.500\ \mathrm{mm} \quad \text{(Dedendum)}
\end{align}

\textbf{OUTSIDE DIAMETER}

\begin{equation}
d_{a,p} = m(N_p+2)
\end{equation}

\begin{equation}
d_{a,p} = 48.000\ \mathrm{mm}
\end{equation}

\begin{equation}
r_{a,p} = 24.000\ \mathrm{mm}
\end{equation}

\textbf{ROOT DIAMETER}

\begin{equation}
d_{f,p} = m(N_p-2.5)
\end{equation}

\begin{equation}
d_{f,p} = 39.000\ \mathrm{mm}
\end{equation}

\begin{equation}
r_{f,p} = 19.500\ \mathrm{mm}
\end{equation}

\textbf{BASE DIAMETER}

\begin{equation}
d_{b,p} = d_p\cos(20^\circ)
\end{equation}

\begin{equation}
d_{b,p} = 41.346\ \mathrm{mm}
\end{equation}

\begin{equation}
r_{b,p} = 20.673\ \mathrm{mm}
\end{equation}

\textbf{CIRCULAR PITCH}

\begin{equation}
p = \pi m
\end{equation}

\begin{equation}
p = 6.283\ \mathrm{mm}
\end{equation}

\textbf{TOOTH THICKNESS}

\begin{equation}
t = \frac{\pi m}{2}
\end{equation}

\begin{equation}
t = 3.1416\ \mathrm{mm}
\end{equation}

\textbf{ANGULAR PITCH}

\begin{equation}
\phi = \frac{360^\circ}{N_p}
\end{equation}

\begin{equation}
\phi = 16.3636^\circ
\end{equation}


\newpage

\subsection{\textbf{Ring Gear ($N_r$ = 64 Teeth)}}

\textbf{PITCH DIAMETER}

\begin{equation}
d_r = mN_r
\end{equation}

\begin{equation}
d_r = 2(64)
\end{equation}

\begin{equation}
d_r = 128.000\ \mathrm{mm}
\end{equation}

\begin{equation}
r_r = 64.000\ \mathrm{mm}
\end{equation}

\textbf{BASE DIAMETER}

\begin{equation}
d_{b,r} = d_r\cos(20^\circ)
\end{equation}

\begin{equation}
d_{b,r} = 120.281\ \mathrm{mm}
\end{equation}

\begin{equation}
r_{b,r} = 60.140\ \mathrm{mm}
\end{equation}

For an internal gear, the tooth geometry is inverted relative to
an external gear.

\textbf{INSIDE TIP / TOOTH-TIP DIAMETER}

\begin{equation}
d_{a,r} = d_r - 2m
\end{equation}

\begin{equation}
d_{a,r} = 128 - 2(2)
\end{equation}

\begin{equation}
d_{a,r} = 124.000\ \mathrm{mm}
\end{equation}

\textbf{OUTSIDE / ROOT DIAMETER}

\begin{equation}
d_{f,r} = d_r + 2.5m
\end{equation}

\begin{equation}
d_{f,r} = 128 + 2.5(2)
\end{equation}

\begin{equation}
d_{f,r} = 133.000\ \mathrm{mm}
\end{equation}


\newpage

\subsection{\textbf{Planet Position}}

\textbf{Centre Distance:}

\begin{align}
C_{s,p} &= \frac{d_s+d_p}{2}
\qquad \text{[Sun--Planet]}\\
&= \frac{40+44}{2}\\
C_{s,p} &= 42.000\ \mathrm{mm}
\end{align}

\begin{align}
C_{p,r} &= \frac{d_r-d_p}{2}
\qquad \text{[Planet--Ring]}\\
&= \frac{128-44}{2}\\
C_{p,r} &= 42.000\ \mathrm{mm}
\end{align}

\textbf{Carrier Centre Radius:}

\begin{align}
R_c &= 42.000\ \mathrm{mm}\\
D_c &= 84.000\ \mathrm{mm}
\end{align}


\subsection{\textbf{Gear Ratio: Sun--Ring}}

With the carrier fixed:

\begin{equation}
i = \frac{N_r}{N_s}
\end{equation}

\begin{equation}
i = \frac{64}{20}
\end{equation}

\begin{equation}
\textbf{i = 3.2:1}
\end{equation}

Using the Willis equation:

\begin{equation}
\frac{\omega_s-\omega_c}{\omega_r-\omega_c}
=
-\frac{N_r}{N_s}
\end{equation}

For a fixed carrier, $\omega_c=0$:

\begin{equation}
\frac{\omega_s}{\omega_r}
=
-\frac{N_r}{N_s}
\end{equation}

Therefore,

\textbf{Given:}

\begin{equation}
\omega_s = 10\ \mathrm{RPM}
\end{equation}

Since

\begin{equation}
\omega_r=-\omega_s\frac{N_s}{N_r}
\end{equation}

\begin{equation}
\omega_r=-10\frac{20}{64}
\end{equation}

\begin{equation}
\boxed{\omega_r=-3.125\ \mathrm{RPM}}
\end{equation}
\end{document}